\chapter{2005年北京卷（秋理）}
\section{单选题}
\begin{problem}
设全集 \(U=\R\)，集合 \(M=\{x\mid x>1\}\)，\(P=\{x\mid x^2>1\}\)，则下列关系中正确的是（\(\quad\)）

\choices
    {\(M=P\)}
    {\(P \not\subseteq M\)}
    {\(M \subsetneq P\)}
    {\(\complement_U M \cap P=\emptyset\)}
\answer{C}
\end{problem}

\begin{problem}
“\(m=\frac{1}{2}\)”是“直线 \((m+2)x+3my+1=0\) 与直线 \((m-2)x+(m+2)y-3=0\) 相互垂直”的（\(\quad\)）

\choices
    {充分必要条件}
    {充分而不必要条件}
    {必要而不充分条件}
    {既不充分也不必要条件}
\answer{B}
\end{problem}

\begin{problem}
若 \(\left|\bs{a}\right|=1\)，\(\left|\bs{b}\right|=2\)，\(\bs{c}=\bs{a}+\bs{b}\)，且 \(\bs{c}\perp\bs{a}\)，则向量 \(\bs{a}\) 与 \(\bs{b}\) 的夹角为（\(\quad\)）

\choices
    {\(30^{\circ}\)}
    {\(60^{\circ}\)}
    {\(120^{\circ}\)}
    {\(150^{\circ}\)}
\answer{C}
\end{problem}

\begin{problem}
从原点向圆 \(x^2+y^2-12y+27=0\) 作两条切线，则该圆夹在两条切线间的劣弧长为（\(\quad\)）

\choices
    {\(\pi\)}
    {\(2\pi\)}
    {\(4\pi\)}
    {\(6\pi\)}
\answer{B}
\end{problem}

\begin{problem}
对任意的锐角 \(\alpha\)，\(\beta\)，下列不等关系中正确的是（\(\quad\)）

\choices
    {\(\sin(\alpha+\beta)>\sin\alpha+\sin\beta\)}
    {\(\sin(\alpha+\beta)>\cos\alpha+\cos\beta\)}
    {\(\cos(\alpha+\beta)<\sin\alpha+\sin\beta\)}
    {\(\cos(\alpha+\beta)<\cos\alpha+\cos\beta\)}
\answer{D}
\end{problem}

\begin{problem}
在正四面体 \(P-ABC\) 中，\(D\)，\(E\)，\(F\) 分别是 \(AB\)，\(BC\)，\(CA\) 的中点，下面四个结论中不成立的是（\(\quad\)）

\choices
    {\(BC\parallel\) 平面 \(PDF\)}
    {\(DF\perp\) 平面 \(PAE\)}
    {平面 \(PDF\perp\) 平面 \(ABC\)}
    {平面 \(PAE\perp\) 平面 \(ABC\)}
\answer{C}
\end{problem}

\begin{problem}
北京《财富》全球论坛期间，某高校有 \(14\) 名志愿者参加接待工作，若每天早，中，晚三班，每班 \(4\) 人，每人每天最多值一班，则开幕式当天不同的排班种数为（\(\quad\)）

\choices
    {\(\mathrm C_{14}^{12}\mathrm C_{12}^{4}\mathrm C_{8}^{4}\)}
    {\(\mathrm C_{14}^{12}\mathrm A_{12}^{4}\mathrm A_{8}^{4}\)}
    {\(\dfrac{\mathrm C_{14}^{12}\mathrm C_{12}^{4}\mathrm C_{8}^{4}}{\mathrm A_{3}^{3}}\)}
    {\(\mathrm C_{14}^{12}\mathrm C_{12}^{4}\mathrm C_{8}^{4}\mathrm A_{3}^{3}\)}
\answer{A}
\end{problem}

\begin{problem}
函数 \(f(x)=\dfrac{\sqrt{1-\cos2x}}{\cos x}\)（\(\quad\)）

\choices
    {在 \(\left[0,\frac{\pi}{2}\right),\left(\frac{\pi}{2},\pi\right]\) 上递增，在 \(\left[\pi,\frac{3\pi}{2}\right),\left(\frac{3\pi}{2},2\pi\right]\) 上递减}
    {在 \(\left[0,\frac{\pi}{2}\right),\left(\pi,\frac{3\pi}{2}\right]\) 上递增，在 \(\left[\frac{\pi}{2},\pi\right),\left(\frac{3\pi}{2},2\pi\right]\) 上递减}
    {在 \(\left(\frac{\pi}{2},\pi\right],\left(\frac{3\pi}{2},2\pi\right]\) 上递增，在 \(\left[0,\frac{\pi}{2}\right),\left(\pi,\frac{3\pi}{2}\right]\) 上递减}
    {在 \(\left[0,\frac{3\pi}{2}\right),\left(\frac{3\pi}{2},\pi\right]\) 上递增，在 \(\left[0,\frac{\pi}{2}\right),\left(\frac{\pi}{2},2\pi\right]\) 上递减}
\answer{A}
\end{problem}

\section{填空题}
\begin{problem}
若 \(z_1=a+2\i\)，\(z_2=3-4\i\)，且 \(\frac{z_1}{z_2}\) 为纯虚数，则实数 \(a\) 的值为\fillinblank{}.
\answer{\(\frac{8}{3}\)}
\end{problem}

\begin{problem}
已知 \(\tan\frac{\alpha}{2}=2\)，则 \(\tan\alpha\) 的值为\fillinblank{}，\(\tan\left(\alpha+\frac{\pi}{4}\right)\) 的值为\fillinblank{}.
\answer{\(-\frac{4}{3}\)，\(-\frac{1}{7}\)}
\end{problem}

\begin{problem}
\(\left(x-\dfrac{1}{x}\right)^6\) 的展开式中的常数项是\fillinblank{}.（用数字作答）
\answer{\(-20\)}
\end{problem}

\begin{solution}
展开式的通项为 \(\mathrm C_6^r x^{6-r}\left(-\dfrac{1}{x}\right)^r=(-1)^r\mathrm C_6^r x^{6-2r}\). 要取得常数项，必须有 \(6-2r=0\)，所以 \(r=3\)，常数项为 \((-1)^3\mathrm C_6^3=-20\).
\end{solution}

\begin{problem}
过原点作曲线 \(y=\e^x\) 的切线，则切点的坐标为\fillinblank{}，切线的斜率为\fillinblank{}.
\answer{\((1,\e)\)，\(\e\)}
\end{problem}

\begin{problem}
对于函数 \(f(x)\) 定义域中任意的 \(x_1\)，\(x_2\)（\(x_1\neq x_2\)），有如下结论：

\begin{circlelist}
\item \(f(x_1+x_2)=f(x_1)\cdot f(x_2)\)；
\item \(f(x_1\cdot x_2)=f(x_1)+f(x_2)\)；
\item \(\dfrac{f(x_1)-f(x_2)}{x_1-x_2}>0\)；
\item \(f\left(\dfrac{x_1+x_2}{2}\right)<\dfrac{f(x_1)+f(x_2)}{2}\).
\end{circlelist}

当 \(f(x)=\lg x\) 时，上述结论中正确结论的序号是\fillinblank{}.
\answer{\(\circled{2}\)、\(\circled{3}\)}
\end{problem}

\begin{problem}
已知 \(n\) 次多项式 \(P_n(x)=a_0x^n+a_1x^{n-1}+\cdots+a_{n-1}x+a_n\). 如果在一种算法中，计算 \(x_0^k\)（\(k=2\)，\(3\)，\(4\)，\(\cdots\)，\(n\)）的值需要 \(k-1\) 次乘法，计算 \(P_3(x_0)\) 的值共需要 \(9\) 次运算（\(6\) 次乘法，\(3\) 次加法），那么计算 \(P_{10}(x_0)\) 的值共需要\fillinblank{} 次运算.

下面给出一种减少运算次数的算法：\(P_0(x)=a_0\)，\(P_{k+1}(x)=xP_k(x)+a_{k+1}\)（\(k=0\)，\(1\)，\(2\)，\(\cdots\)，\(n-1\)）. 利用该算法，计算 \(P_3(x_0)\) 的值共需要 \(6\) 次运算，计算 \(P_{10}(x_0)\) 的值共需要\fillinblank{} 次运算.
\answer{\(65\)，\(20\)}
\end{problem}

\section{解答题}
\begin{problem}
已知函数 \(f(x)=-x^3+3x^2+9x+a\).

\begin{enumerate}
\item 求 \(f(x)\) 的单调减区间；
\item 若 \(f(x)\) 在区间 \([-2,2]\) 上的最大值为 \(20\)，求它在该区间上的最小值.
\end{enumerate}
\end{problem}

\begin{solution}
\begin{enumerate}
\item \(f'(x)=-3x^2+6x+9=-3(x+1)(x-3)\). 当 \(x<-1\) 或 \(x>3\) 时，\(f'(x)<0\)；当 \(-1<x<3\) 时，\(f'(x)>0\). 所以 \(f(x)\) 的单调减区间为 \(\left(-\infty,-1\right)\) 和 \(\left(3,+\infty\right)\).

\item 在区间 \(\left[-2,2\right]\) 上，函数在 \(\left[-2,-1\right]\) 上递减，在 \(\left[-1,2\right]\) 上递增. 计算得 \(f(-2)=a+2\)，\(f(-1)=a-5\)，\(f(2)=a+22\).
其中 \(f(2)>f(-2)\)，故区间上的最大值为 \(f(2)=a+22=20\)，从而 \(a=-2\). 因此最小值为 \(f(-1)=a-5=-7\).
\end{enumerate}
\end{solution}

\begin{problem}
如图，在直四棱柱 \(ABCD-A_{1}B_{1}C_{1}D_{1}\) 中，\(AB=AD=2\)，\(DC=2\sqrt{3}\)，\(AA_1=\sqrt{3}\)，\(AD\perp DC\)，\(AC\perp BD\)，垂足为 \(E\).

\begin{enumerate}
\item 求证：\(BD\perp A_1C\)；
\item 求二面角 \(A_1-BD-C_1\) 的大小；
\item 求异面直线 \(AD\) 与 \(BC_1\) 所成角的大小.
\end{enumerate}

\bitmapfigure[max width=0.45\linewidth]{img/2005/beijing_science_autumn/q16_fig1.png}
\end{problem}

\begin{solution}
建立空间直角坐标系，取 \(D\) 为原点，\(DC\) 所在直线为 \(x\) 轴，\(DA\) 所在直线为 \(y\) 轴，直四棱柱的侧棱方向为 \(z\) 轴. 于是 \(D=(0,0,0)\)，\(A=(0,2,0)\)，\(C=(2\sqrt{3},0,0)\). 设 \(B=(u,v,0)\). 由 \(AB=2\) 及 \(AC\perp BD\)，分别得到
\[
u^2+(v-2)^2=4.
\]
\[
(2\sqrt{3},-2,0)\cdot(-u,-v,0)=0.
\]
所以 \(v=\sqrt{3}u\). 代入前式得 \(u^2+(\sqrt{3}u-2)^2=4\)，即 \(4u(u-\sqrt{3})=0\). 由于四棱柱的底面不退化，\(B\ne D\)，故 \(u=\sqrt{3}\)，\(v=3\)，从而 \(B=(\sqrt{3},3,0)\)，\(A_1=(0,2,\sqrt{3})\)，\(C_1=(2\sqrt{3},0,\sqrt{3})\).

\begin{enumerate}
\item \(\overrightarrow{BD}=(-\sqrt{3},-3,0)\)，\(\overrightarrow{A_1C}=(2\sqrt{3},-2,-\sqrt{3})\)，所以
\[
\overrightarrow{BD}\cdot\overrightarrow{A_1C}=-6+6=0.
\]
因此 \(BD\perp A_1C\).

\item 平面 \(A_1BD\) 的两个方向向量可取 \(\overrightarrow{DB}=(\sqrt{3},3,0)\) 和 \(\overrightarrow{DA_1}=(0,2,\sqrt{3})\)，平面 \(C_1BD\) 的两个方向向量可取 \(\overrightarrow{DB}\) 和 \(\overrightarrow{DC_1}=(2\sqrt{3},0,\sqrt{3})\). 因此这两个平面的一个法向量分别为 \((3\sqrt{3},-3,2\sqrt{3})\) 和 \((3\sqrt{3},-3,-6\sqrt{3})\).
它们的数量积为 \(27+9-36=0\)，故两平面垂直，二面角 \(A_1-BD-C_1\) 的大小为 \(90^\circ\).

\item 取 \(\overrightarrow{AD}=(0,-2,0)\)，则 \(\overrightarrow{BC_1}=(\sqrt{3},-3,\sqrt{3})\). 因此
\(\left|\overrightarrow{AD}\cdot\overrightarrow{BC_1}\right|=6\)，\(\left|\overrightarrow{AD}\right|=2\)，\(\left|\overrightarrow{BC_1}\right|=\sqrt{15}\).
设异面直线所成的角为 \(\theta\)，则 \(\cos\theta=\dfrac{6}{2\sqrt{15}}=\dfrac{\sqrt{15}}{5}\)，所以 \(\theta=\arccos\dfrac{\sqrt{15}}{5}\).
\end{enumerate}
\end{solution}

\begin{problem}
甲，乙两人各进行 \(3\) 次射击，甲每次击中目标的概率为 \(\frac{1}{2}\)，乙每次击中目标的概率为 \(\frac{2}{3}\).

\begin{enumerate}
\item 记甲击中目标的次数为 \(\xi\)，求 \(\xi\) 的概率分布及数学期望 \(E(\xi)\)；
\item 求乙至多击中目标 \(2\) 次的概率；
\item 求甲恰好比乙多击中目标 \(2\) 次的概率.
\end{enumerate}
\end{problem}

\begin{solution}
设乙击中目标的次数为 \(\eta\). 由于各次射击相互独立，\(\xi\) 服从参数为 \(3\)，\(\frac{1}{2}\) 的二项分布，故
\begin{enumerate}
\item 对 \(k=0,1,2,3\)，有 \(P(\xi=k)=\mathrm C_3^k\left(\dfrac{1}{2}\right)^k\left(\dfrac{1}{2}\right)^{3-k}=\dfrac{\mathrm C_3^k}{8}\). 因此
\(P(\xi=0)=\frac18\)，\(P(\xi=1)=\frac38\)，\(P(\xi=2)=\frac38\)，\(P(\xi=3)=\frac18\).
并且
\[
E(\xi)=0\cdot\frac18+1\cdot\frac38+2\cdot\frac38+3\cdot\frac18=\frac32.
\]

\item \(\eta\) 服从参数为 \(3\)，\(\frac23\) 的二项分布，所以
\[
P(\eta\leqslant2)=1-P(\eta=3)=1-\left(\frac23\right)^3=\frac{19}{27}.
\]

\item 甲恰好比乙多击中目标 \(2\) 次时，可能的情形为 \((\xi,\eta)=(2,0)\) 或 \((3,1)\). 由于甲、乙的射击相互独立，所求概率为
\[
P(\xi=2)P(\eta=0)+P(\xi=3)P(\eta=1)
=\frac38\cdot\frac1{27}+\frac18\cdot\mathrm C_3^1\frac23\left(\frac13\right)^2
=\frac1{72}+\frac1{36}=\frac1{24}.
\]
\end{enumerate}
\end{solution}

\begin{problem}
如图，直线 \(l_1:y=kx\)（\(k>0\)）与直线 \(l_2:y=-kx\) 之间的阴影区域（不含边界）记为 \(W\)，其左半部分记为 \(W_1\)，右半部分记为 \(W_2\).

\begin{enumerate}
\item 分别用不等式组表示 \(W_1\) 和 \(W_2\)；
\item 若区域 \(W\) 中的动点 \(P(x,y)\) 到 \(l_1\)，\(l_2\) 的距离之积等于 \(d^2\)，求点 \(P\) 的轨迹 \(C\) 的方程；
\item 设不过原点 \(O\) 的直线 \(l\) 与（2）中的曲线 \(C\) 相交于 \(M_1\)，\(M_2\) 两点，且与 \(l_1\)，\(l_2\) 分别交于 \(M_3\)，\(M_4\) 两点. 求证：\(\triangle OM_1M_2\) 的重心与 \(\triangle OM_3M_4\) 的重心重合.
\end{enumerate}

\bitmapfigure[max width=0.42\linewidth]{img/2005/beijing_science_autumn/q18_fig1.png}
\end{problem}

\begin{solution}
\begin{enumerate}
\item 当 \(x<0\) 时，\(kx<-kx\)，故左半部分为
\[
W_1=\{(x,y)\mid x<0,\ kx<y<-kx\}.
\]
当 \(x>0\) 时，\(-kx<kx\)，故右半部分为
\[
W_2=\{(x,y)\mid x>0,\ -kx<y<kx\}.
\]

\item 直线 \(l_1\)，\(l_2\) 的方程分别为 \(kx-y=0\)，\(kx+y=0\)，所以点 \(P(x,y)\) 到它们的距离分别为 \(\frac{|kx-y|}{\sqrt{1+k^2}}\) 和 \(\frac{|kx+y|}{\sqrt{1+k^2}}\).
在区域 \(W\) 内有 \(y^2<k^2x^2\)，因此距离之积为
\[
\frac{|kx-y||kx+y|}{1+k^2}=\frac{k^2x^2-y^2}{1+k^2}.
\]
由题意，该积等于 \(d^2\)，故轨迹 \(C\) 的方程为
\[
k^2x^2-y^2-(1+k^2)d^2=0.
\]

\item 设不过原点的直线 \(l\) 为 \(y=mx+n\). 由于 \(O\notin l\)，有 \(n\ne0\). 设 \(M_1\)，\(M_2\) 的横坐标为 \(x_1\)，\(x_2\)，代入曲线方程得
\[
(k^2-m^2)x^2-2mnx-n^2-(1+k^2)d^2=0.
\]
因该直线与曲线有两个交点，\(m^2\ne k^2\)，由韦达定理
\[
x_1+x_2=\frac{2mn}{k^2-m^2}.
\]
又设 \(M_3\)，\(M_4\) 的横坐标为 \(x_3\)，\(x_4\)，则 \(x_3=\frac{n}{k-m}\)，\(x_4=-\frac{n}{k+m}\)，从而 \(x_3+x_4=\frac{2mn}{k^2-m^2}\).
所以 \(x_1+x_2=x_3+x_4\). 由于四点均在直线 \(l\) 上，纵坐标满足 \(y_i=mx_i+n\)，从而
\[
y_1+y_2=m(x_1+x_2)+2n=m(x_3+x_4)+2n=y_3+y_4.
\]
因此
\[
\overrightarrow{OM_1}+\overrightarrow{OM_2}=\overrightarrow{OM_3}+\overrightarrow{OM_4}.
\]
三角形 \(OM_1M_2\) 与 \(OM_3M_4\) 的重心坐标分别为上述两式的三分之一，故两三角形的重心重合.

若 \(l\) 为竖直直线 \(x=c\)，则 \(c\ne0\)，且 \(M_3=(c,kc)\)，\(M_4=(c,-kc)\)，两点坐标和为 \((2c,0)\)；曲线与该直线的两个交点的纵坐标互为相反数，坐标和同样为 \((2c,0)\)，结论仍成立.
\end{enumerate}
\end{solution}

\begin{problem}
设数列 \(\{a_n\}\) 的首项 \(a_1=a\neq\frac{1}{4}\)，且
\(a_{n+1}=
\begin{cases}
\dfrac{1}{2}a_n, & n\text{ 为偶数},\\
a_n+\dfrac{1}{4}, & n\text{ 为奇数}.
\end{cases}\)
记 \(b_n=a_{2n-1}-\dfrac{1}{4}\)，\(n=1\)，\(2\)，\(3\)，\(\cdots\).

\begin{enumerate}
\item 求 \(a_2\)，\(a_3\)；
\item 判断数列 \(\{b_n\}\) 是否为等比数列，并证明你的结论；
\item 求 \(\displaystyle\lim_{n\to\infty}\left(b_1+b_2+\cdots+b_n\right)\).
\end{enumerate}
\end{problem}

\begin{solution}
\begin{enumerate}
\item 因为 \(1\) 为奇数，\(a_2=a_1+\frac14=a+\frac14\). 因为 \(2\) 为偶数，\(a_3=\frac12a_2=\frac12a+\frac18\).

\item \(b_1=a_1-\frac14=a-\frac14\). 对任意正整数 \(n\)，由递推关系有 \(a_{2n}=a_{2n-1}+\frac14\)，\(a_{2n+1}=\frac12a_{2n}\).
所以
\[
b_{n+1}=a_{2n+1}-\frac14=\frac12\left(a_{2n-1}+\frac14\right)-\frac14=\frac12\left(a_{2n-1}-\frac14\right)=\frac12b_n.
\]
又因为 \(a\ne\frac14\)，所以 \(b_1\ne0\)，从而 \(\{b_n\}\) 是首项为 \(a-\frac14\)，公比为 \(\frac12\) 的等比数列.

\item 由等比数列求和公式
\[
b_1+b_2+\cdots+b_n=\left(a-\frac14\right)\frac{1-\left(\frac12\right)^n}{1-\frac12}.
\]
因为 \(\left(\frac12\right)^n\to0\)，所以
\[
\lim_{n\to\infty}(b_1+b_2+\cdots+b_n)=2\left(a-\frac14\right)=2a-\frac12.
\]
\end{enumerate}
\end{solution}

\begin{problem}
设 \(f(x)\) 是定义在 \([0,1]\) 上的函数，若存在 \(x^{*}\in(0,1)\)，使得 \(f(x)\) 在 \([0,x^{*}]\) 上单调递增，在 \([x^{*},1]\) 上单调递减，则称 \(f(x)\) 为 \([0,1]\) 上的单峰函数，\(x^{*}\) 为峰点，包含峰点的区间为含峰区间. 对任意的 \([0,1]\) 上的单峰函数 \(f(x)\)，下面研究缩短其含峰区间长度的方法.

\begin{enumerate}
\item 证明：对任意的 \(x_1\)，\(x_2\in(0,1)\)，\(x_1<x_2\)，\(f(x_1)\geqslant f(x_2)\)，则 \(\left(0,x_2\right)\) 为含峰区间；若 \(f(x_1)\leqslant f(x_2)\)，则 \(\left(x_1,1\right)\) 为含峰区间；
\item 对给定的 \(r\)（\(0<r<0.5\)），证明：存在 \(x_1\)，\(x_2\in(0,1)\)，满足 \(x_2-x_1\geqslant2r\)，使得由（1）所确定的含峰区间的长度不大于 \(0.5+r\)；
\item 选取 \(x_1\)，\(x_2\in(0,1)\)，\(x_1<x_2\)，由（1）可确定含峰区间为 \(\left(0,x_2\right)\) 或 \(\left(x_1,1\right)\)，在所得的含峰区间内选取 \(x_3\)，由 \(x_3\) 与 \(x_1\) 或 \(x_3\) 与 \(x_2\) 类似地可确定一个新的含峰区间，在第一次确定的含峰区间为 \(\left(0,x_2\right)\) 的情况下，试确定 \(x_1\)，\(x_2\)，\(x_3\) 的值，满足两两之差的绝对值不小于 \(0.02\)，且使得新的含峰区间的长度缩短到 \(0.34\).
\end{enumerate}

注：区间长度等于区间的右端点与左端点之差.
\end{problem}

\begin{solution}
\begin{enumerate}
\item 设峰点为 \(x^*\). 若 \(f(x_1)\geqslant f(x_2)\)，则峰点不能在 \(x_2\) 的右侧；否则 \(x_1\)，\(x_2\) 均处于递增段，将有 \(f(x_1)<f(x_2)\). 因此可取 \(\left(0,x_2\right)\) 为含峰区间. 类似地，若 \(f(x_1)\leqslant f(x_2)\)，则峰点不能在 \(x_1\) 的左侧；否则 \(x_1\)，\(x_2\) 均处于递减段，将有 \(f(x_1)>f(x_2)\). 因此可取 \(\left(x_1,1\right)\) 为含峰区间.

\item 取 \(x_1=0.5-r\)，\(x_2=0.5+r\). 由于 \(0<r<0.5\)，有 \(x_1\in(0,1)\)，\(x_2\in(0,1)\)，且 \(x_2-x_1=2r\). 由第（1）问确定的含峰区间只有两种可能：\(\left(0,x_2\right)\) 或 \(\left(x_1,1\right)\). 这两个区间的长度分别为 \(x_2=0.5+r\) 和 \(1-x_1=0.5+r\).
故所确定的含峰区间长度不大于 \(0.5+r\).

\item 取 \(x_1=0.34\)，\(x_2=0.66\)，\(x_3=0.32\). 三点均在 \((0,1)\) 内，且 \(|x_2-x_1|=0.32\)，\(|x_1-x_3|=0.02\)，\(|x_2-x_3|=0.34\)，满足两两之差的绝对值不小于 \(0.02\). 在第一次确定的含峰区间为 \(\left(0,x_2\right)=\left(0,0.66\right)\) 的条件下，将 \(x_3\) 与 \(x_1\) 比较：

若 \(f(x_3)\geqslant f(x_1)\)，由第（1）问可取新的含峰区间为 \(\left(0,x_1\right)=\left(0,0.34\right)\)，其长度为 \(0.34\)；若 \(f(x_3)\leqslant f(x_1)\)，由第（1）问得到 \(\left(x_3,1\right)\)，与原含峰区间取交集后得到 \(\left(x_3,x_2\right)=\left(0.32,0.66\right)\)，其长度仍为 \(0.34\). 因此所取三点满足要求.
\end{enumerate}
\end{solution}
